% ===========================================================================
%  Macros for the Trainingsblaetter.
%
%  Self-contained on purpose: the worksheets have their own title page
%  (\HAtitle) and their own task/solution layout, and share only the package
%  list with the rest of the repository (gemeinsam/preamble/packages.tex).
% ===========================================================================

\newcommand{\HAtitle}[2]{
	\begin{titlepage} % Suppresses displaying the page number on the title page and the subsequent page counts as page 1
		\newcommand{\HRule}{\rule{\linewidth}{0.5mm}} % Defines a new command for horizontal lines, change thickness here
		\center % Centre everything on the page
		\hspace{1cm} \\[3.5cm]
		\HRule\\[0.4cm]
		{\huge\bfseries Numerisches Programmieren Trainingsaufgaben #1}\\[0.4cm] % Title of your document
		\textsc{\Large #2}\\[0.5cm] % Minor heading such as course title
		\HRule\\[1.5cm]
		\begin{minipage}{0.4\textwidth}
			\begin{flushleft}
				\large
				\textit{Autor}\\
				\textsc{\npauthor{}} % Your name
			\end{flushleft}
		\end{minipage}
		~
		\begin{minipage}{0.4\textwidth}
			\begin{flushright}
				\large
				\textit{Version vom:}\\
				\textsc{\today} % Supervisor's name
			\end{flushright}
		\end{minipage}
		\vfill\vfill % Position the date 3/4 down the remaining page
		\large{\textit{Erst die Aufgaben rechnen, bevor Ihr die Lösung anguckt!}}
		\vfill % Push the date up 1/4 of the remaining page
	\end{titlepage}
}

\definecolor{easy}{rgb}{0.2 0.2 1}
\definecolor{medium}{rgb}{0.9 0.4 0}
\definecolor{hard}{rgb}{1 0 0.4}
\definecolor{special}{rgb}{0.5 0.2 1}
\newcommand{\easy}{\textcolor{easy}{(Easy)}}
\newcommand{\medium}{\textcolor{medium}{(Medium)}}
\newcommand{\hard}{\textcolor{hard}{(Hard)}}
\newcommand{\transfer}{\textcolor{special}{(Transfer)}}

\newcommand{\itemf}{\item}
\newcommand{\itemeasy}{\item[\color{easy}]\alph*)}


\newcommand{\task}[4]{
	\clearpage
	\section{#1}
	%
	#2
	\begin{enumerate}[label=\alph*)]
		#3
	\end{enumerate}
	
	\vfill \textit{Auf der nächsten Seite folgt die Lösung\dots}
	\clearpage
	\Large \textbf{Lösung:}
	\vspace{0.3cm}
	\normalsize
	\begin{enumerate}[label=\alph*), nolistsep, itemsep=\fill]
		#4
	\end{enumerate}
	\vfill
}

\newcommand{\subtask}[1]{
	\begin{enumerate}[label=(\roman*)]
		#1
	\end{enumerate}
}

\newcommand{\hinweis}{\textit{Hinweis: }}
% =========================================================
% Specific math variables:
% =========================================================

\newcommand{\fabs}{f_{\text{abs}}}
\newcommand{\frel}{f_{\text{rel}}}

\newcommand{\maxError}{\epsilon_{\text{Ma}}}

\newcommand{\mddots}{\reflectbox{$\ddots$}}


\newcommand{\umax}{_\text{max}}
\newcommand{\umin}{_\text{min}}

% =========================================================
% Math Commands
% =========================================================

% Should be renamed or removed
\newcommand{\st}[1]{^{(#1)}}

% vertical line in the middle:
\newcommand{\separator}[1]{\; #1 \;}
\newcommand{\sep}{\separator{|}}
\newcommand{\sepArrow}{\separator{\Longleftrightarrow}}

% BigO
\newcommand{\bigO}{\mathcal{O}} % Big-O notation

% Norm with double |

% Matrix Transpose and inverted
\newcommand{\trans}{^{\mathrm{\mathsmaller T}}}
\newcommand{\inv}{^{\mathrm{\mathsmaller{\text{-}1}}}}
\newcommand{\unitM}{\mathbb{1}}


% informal descriptions:
\newcommand{\uleftDown}{_\text{linksUnten}}
\newcommand{\uleftUp}{_\text{linksOben}}
\newcommand{\uleft}{_\text{links}}
\newcommand{\uleftDiagonal}{_\text{ganzLinksUnten}}
\newcommand{\uleftUpDiagonal}{_\text{ganzLinksOben}}

% Soll sein
\newcommand{\shouldEq}{ \stackrel{!}{=} }

% Formelnamen:
\newcommand{\QT}{Q_{\text{T}}}
\newcommand{\QTS}{Q_{\text{TS}}}
\newcommand{\QS}{Q_{\text{S}}}
\newcommand{\QSS}{Q_{\text{SS}}}
\newcommand{\RTS}{R_{\text{TS}}}
\newcommand{\RSS}{R_{\text{SS}}}
%--------------------------

% correct strange spacing around delimiters
\renewcommand{\l}{\mathopen{}\mathclose\bgroup\left}
\renewcommand{\r}{\aftergroup\egroup\right}

% norms
\newcommand{\norm}[1]{\l\lVert #1 \r\rVert}
\newcommand{\pnorm}[2]{\norm{#1}_{#2}}

\newcommand{\nvector}[1]{
	\begin{pmatrix}
		#1
	\end{pmatrix}
}

\newcommand{\nmatrix}[1]{
	\begin{bmatrix}
		#1
	\end{bmatrix}
}

\def\nLongrightarrow{\centernot \Longrightarrow}
\def\nLongleftrightarrow{\centernot \Longleftrightarrow}
\def\nforall{\centernot \forall}
\def\nexists{\centernot \exists}
\def\nsqsubseteq{\centernot \sqsubseteq}
\def\nsubset{\centernot \subset}
\def\nequiv{\centernot \equiv}

\def\symdif{\triangle}
\def\imp{\rightarrow}
\def\bik{\leftrightarrow}
\def\xor{\oplus}
\def\nand{\, \bar{\land}\,}
\def\nor{\,\bar{\lor}\,}


% number sets
\newcommand{\booleans}[0]{\mathbb{B}}			% Perhaps call it \binaries and defined it as \residuals{2} instead.
\newcommand{\naturals}[0]{\mathbb{N}}
\newcommand{\wholes}[0]{\mathbb{N}_0}
\newcommand{\integers}[0]{\mathbb{Z}}
\newcommand{\rationals}[0]{\mathbb{Q}}
\newcommand{\reals}[0]{\mathbb{R}}
\newcommand{\complexes}[0]{\mathbb{C}}
\newcommand{\residuals}[1]{\mathbb{Z}_{#1}}
\newcommand{\primes}[0]{\mathbb{P}}

% brackets (round, square, curly, angled)
\renewcommand{\(}[0]{\l(}						% \(...\) was defined for $...$ as \[...\] for $$...$$
\renewcommand{\)}[0]{\r)}
\newcommand{\rb}[1]{\l(#1\r)}
\renewcommand{\sb}[1]{\l[#1\r]}
\newcommand{\cb}[1]{\l\{#1\r\}}
\newcommand{\ab}[1]{\l\langle#1\r\rangle}

% generating set
\newcommand{\generatingset}[1]{\ab{#1}}

% floor and ceiling
\newcommand{\ceil}[1]{\l\lceil#1\r\rceil}
\newcommand{\floor}[1]{\l\lfloor#1\r\rfloor}


% absolute value
\newcommand{\abs}[1]{\l|#1\r|}

% integrals
\newcommand{\integral}[4]{\int_{#1}^{#2} #3 \ \mathrm{d}#4}
\newcommand{\antiderivative}[2]{\integral{}{}{#1}{#2}}
\renewcommand{\d}[1]{\ \mathrm{d}#1}

% logarithms
\newcommand{\ld}[0]{\operatorname{ld}}
\newcommand{\lb}[0]{\operatorname{lb}}

% signum function
\newcommand{\sgn}[0]{\operatorname{sgn}}

% fixed set
\newcommand{\fix}[0]{\operatorname{fix}}		% fix(f) = {x|f(x) = x}

% closures
\newcommand{\refl}[0]{\operatorname{ref}}
\newcommand{\symm}[0]{\operatorname{sym}}

% sets of functions
\newcommand{\inj}[0]{\operatorname{inj}}
\newcommand{\surj}[0]{\operatorname{surj}}
\newcommand{\bij}[0]{\operatorname{bij}}

% fields
\newcommand{\charac}[0]{\operatorname{char}}
\newcommand{\GF}[0]{\operatorname{GF}}

% least common multiple and greatest common divisor
\newcommand{\ggT}[0]{\operatorname{ggT}}
\newcommand{\kgV}[0]{\operatorname{kgV}}

% average
\newcommand{\avg}[0]{\operatorname{avg}}

% order
\newcommand{\ord}[0]{\operatorname{ord}}

% ----------
% round 
\newcommand{\rd}[0]{\operatorname{rd}}
\newcommand{\rdM}[0]{\operatorname{rd_M}}

% condition value
\newcommand{\cond}[0]{\operatorname{cond}}

% complex numbers
\newcommand{\realPart}{\operatorname{Re}}
\newcommand{\imaginaryPart}{\operatorname{Im}}

% Fourier transformation
\newcommand{\DFT}{\operatorname{DFT}}
\newcommand{\IDFT}{\operatorname{IDFT}}
\newcommand{\FFT}{\operatorname{FFT}}
\newcommand{\BFO}{\operatorname{BFO}}

% diagonal entries of a matrix
\newcommand{\diag}{\operatorname{diag}}
% --------

% sequent calculus
\newcommand{\FL}[0]{\operatorname{FL}}
\newcommand{\Nat}[0]{\operatorname{Nat}}

% true and false
\newcommand{\false}[0]{\operatorname{false}}
\newcommand{\true}[0]{\operatorname{true}}

% expected value
\newcommand{\Eempty}[0]{\mathbb{E}}
\newcommand{\E}[1]{\Eempty\sb{#1}}
\newcommand{\Eindex}[2]{\Eempty_{#1}\sb{#2}}
\newcommand{\Econd}[2]{\Eempty\sb{#1 \middle| #2}}	% E[#1|#2]

% variance
\newcommand{\Varempty}[0]{\operatorname{Var}}
\newcommand{\Var}[1]{\Varempty\sb{#1}}
\newcommand{\Varindex}[2]{\Varempty_{#1}\sb{#2}}
\newcommand{\Varcond}[2]{\Varempty\sb{#1 \middle| #2}}	% Var[#1|#2]

\newcommand{\MSE}[0]{\operatorname{MSE}}
\newcommand{\Bias}[0]{\operatorname{Bias}}

% probability distributions
\newcommand{\Ber}[0]{\operatorname{Ber}}
\newcommand{\Bin}[0]{\operatorname{Bin}}
\newcommand{\Poi}[0]{\operatorname{Poi}}
\newcommand{\Geo}[0]{\operatorname{Geo}}
\newcommand{\NegBin}[0]{\operatorname{NegBin}}
\newcommand{\Hyp}[0]{\operatorname{Hyp}}
\newcommand{\NegHyp}[0]{\operatorname{NegHyp}}
\newcommand{\Nor}[0]{\operatorname{\mathcal{N}}}
\newcommand{\Exp}[0]{\operatorname{Exp}}
\newcommand{\Erl}[0]{\operatorname{Erl}}
\newcommand{\ConUni}[0]{\operatorname{Uni}}
\newcommand{\DisUni}[0]{\operatorname{Uni}}
\newcommand{\Uni}[0]{\operatorname{Uni}}					% TODO: cambiar por \UniDis y \UniCon

% relations
\newcommand{\id}[1]{\operatorname{id}_{#1}}					% Identitätsrelation
\newcommand{\all}[1]{#1 \times #1}							% Allrelation

% =========================================================
% TikZ Commands
% =========================================================

\newcommand{\tikzPoint}[3]{
	\node[label={#1:{#2}},circle,fill,inner sep=2pt] at (axis cs:#3) {};
}
\newcommand{\tikzPointLeft}[2]{
	\tikzPoint{180}{#1}{#2}
}
\newcommand{\tikzPointRight}[2]{
	\tikzPoint{0}{#1}{#2}
}
\newcommand{\tikzPointUp}[2]{
	\tikzPoint{90}{#1}{#2}
}
\newcommand{\tikzPointDown}[2]{
	\tikzPoint{270}{#1}{#2}
}
\newcommand{\tikzPointUL}[2]{
	\tikzPoint{135}{#1}{#2}
}
\newcommand{\tikzPointUR}[2]{
	\tikzPoint{45}{#1}{#2}
}
\newcommand{\tikzPointDL}[2]{
	\tikzPoint{225}{#1}{#2}
}
\newcommand{\tikzPointDR}[2]{
	\tikzPoint{315}{#1}{#2}
}

%[red, very thick, |-] (0,0) -- (22,0);
\def \tikzPattern{on 1pt off 3pt on 3pt off 3pt}%dash pattern=on 1pt off 3pt on 3pt off 3pt}

\newcommand{\tikzLineDotted}[3]{
	\draw[gray, dotted] (axis cs:#2) -- node[left]{#1} (axis cs:#3);
}
\newcommand{\tikzLine}[4]{
	\draw[#1] (axis cs:#3) -- node[left]{#2} (axis cs:#4);
}
\newcommand{\tikzLineUp}[4]{
	\draw[#1] (axis cs:#3) -- node[above]{#2} (axis cs:#4);
}
\newcommand{\tikzNormalLine}[3]{
	\tikzLine{black, thick}{#1}{#2}{#3}
}
\newcommand{\tikzFatLine}[3]{
	\tikzLine{blue, very thick}{#1}{#2}{#3}
}
\newcommand{\tikzFatterLine}[3]{
	\tikzLine{blue, line width=0.65mm}{#1}{#2}{#3}
}
\newcommand{\tikzMegaFatLine}[3]{
	\tikzLine{blue, line width=0.95mm}{#1}{#2}{#3}
}

\newcommand{\tikzArrow}[3]{
	\draw[very thick, ->] (axis cs:#2) -- node[above]{#1} (axis cs:#3);
}
% Angles:
% 0 -> right, 90 = up, 270 = down, 180
\newcommand{\tikzCircle}[3]{
	\draw(axis cs:#2) circle[#1, radius=#3];
}

\newcommand{\tikzTrigonometric}{
	xtick={-3.14159, -1.5708, 0,1.5708,3.14159,4.7123889},
	xticklabels={$-\pi$, $-\frac{\pi}{2}$, $0$,$\frac{\pi}{2}$,$\pi$,$\frac{3\pi}{2}$},
}


\newcommand{\tikzFigure}[2]{
	\begin{figure}[ht]
		\centering
		\begin{tikzpicture}
		#1
		\end{tikzpicture}
		\caption{#2}
	\end{figure}
}
\newcommand{\tikzFigureStandard}[4]{
	\tikzFigure{
		\begin{axis}[
			axis lines=left,
			xlabel=$x$,
			ylabel=$y$,
			enlargelimits,
			%xtick={-3.14159, -1.5708, 0,1.5708,3.14159,4.7123889},
			%xticklabels={$-\pi$, $-\frac{\pi}{2}$, $0$,$\frac{\pi}{2}$,$\pi$,$\frac{3\pi}{2}$},
			%ymax=2,
			%ymin=-0.5]
			%[domain=0:1,legend pos=outer north east]
			\addplot[domain=#1:#2] {#3} node[left]{};
			%				( 0, 4 / pi )
			%				( pi, 0 )
			%				( 0, 0)
			%				( 0, 4 / pi)
			%			};
		\end{axis}
	}{#4}
}